% -----------------------------------------------
% DJ Metadata Paper — Full Draft with New Sections
% Paste into main.tex replacing corresponding sections
% -----------------------------------------------

% ---- RELATED WORK (expanded) ----

\section{Related Work}

\textbf{Music metadata interoperability.} The MusicBrainz project~\citep{musicbrainz} and Discogs~\citep{discogs} provide community-standard schemas for general music metadata but do not capture DJ-specific fields such as BPM, musical key notation, energy/mood ratings, or playlist structure. Hardjono et al.~\citep{hardjono2019} propose an open music metadata layer using distributed ledgers but do not analyze DJ-platform interoperability.

The Music Meta Ontology~\citep{deberardinis2023metontology} proposes a flexible semantic model for music metadata interoperability across domains---directly relevant to the cross-platform DJ metadata problem. Discogs-VI~\citep{araz2024discogsvi} demonstrates the value of editorial metadata from Discogs for version identification, further motivating the need for standardized cross-platform metadata schemas. Scherpenisse's thesis~\citep{scherpenisse2004} studies music metadata management in the context of MusicBrainz, including metadata survey, data cleaning, and distributed server architectures, but focuses on music identification rather than cross-platform library transfer. Correya et al.~\citep{correya2018coversong} show that combining metadata with audio features improves cover song detection at scale, demonstrating metadata's utility for identification tasks.

\textbf{DJ software ecosystems and mix analysis.} Rekordbox XML exports have been reverse-engineered by the community~\citep{pyrekordbox}. Serato's database format has been analyzed via the Serato-Library-Parser project~\citep{serato-parser}. Traktor's NML format remains partially documented. No prior work has conducted a systematic cross-platform transfer study in the DJ domain. However, recent work on DJ mix analysis provides context: Kim et al.~\citep{kim2020djmixes} analyze real-world DJ mixes using subsequence alignment to extract track transitions, while Andr{\'e} et al.~\citep{andre2024djmix} reverse-engineer DJ mixes via non-negative matrix factorization to recover source tracks and applied effects. Zehren et al.~\citep{zehren2020cuepoints} address automatic cue point detection for DJ mixing, a DJ-specific metadata field. These works demonstrate that DJ metadata---cue points, transitions, mix structure---is analyzable at scale but depends on consistent cross-platform representation.

\textbf{Beat tracking and key detection.} BPM is the most robustly preserved field in our study, consistent with the maturity of beat tracking systems. BeatNet~\citep{heydari2021beatnet} achieves online beat and downbeat tracking via CRNN and particle filtering, while BEAST~\citep{chang2024beast} extends this to streaming transformers for real-time applications. Chiu et al.~\citep{chiu2022beatswitching} analyze metric-level switching in beat tracking, highlighting cases where tempo metadata may be ambiguous---relevant to edge cases in cross-platform transfer. Silva~\citep{silva2025keydetection} provides a detailed analysis of audio key detection algorithms, explaining why key notation varies across platforms (OpenKey, Camelot, traditional).

\textbf{Audio fingerprinting and version identification.} Ke et al.~\citep{ke2005audio} establish computer vision-based audio fingerprinting for music identification, while Kurth and M{\"u}ller~\citep{kurth2008efficient} demonstrate scalable index-based audio matching. Jang et al.~\citep{jang2009pairwise} apply pairwise boosting to improve fingerprint robustness. Yu et al.~\citep{yu2019coversong} learn CNN representations for cover song identification. These techniques form the basis for auto-tag recovery when metadata degrades during transfer---a key component of our Experiment~5.

\textbf{Music tagging and ML metadata dependence.} TTMR++~\citep{ttmr2024} enriches music descriptions using LLMs and metadata for retrieval, demonstrating that structured metadata improves downstream tasks. Choi et al.~\citep{choi2019zeroshot} show zero-shot music classification depends on knowledge transfer from existing metadata. Bukey et al.~\citep{bukey2026musicmetadata} rethink music captioning with metadata-aware LLMs. These works collectively establish that ML systems for music---including DJ-specific tools like CUE-DETR~\citep{arguello2024cuedetr} for cue point detection---depend on metadata quality. Our work quantifies how much metadata is lost in practice and provides a schema to prevent it.

Doh et al.~\citep{doh2021playlist} generate playlist titles via machine translation, while Liebman et al.~\citep{liebman2015djmc} use reinforcement learning for DJ playlist recommendation. Both depend on structured metadata for playlist representation---a field particularly vulnerable to cross-platform degradation.


% ---- METHODS (expanded) ----

\section{Methods}

\subsection{Dataset}

We assembled a dataset of 2,147 tracks spanning electronic (41\%), hip-hop (22\%), rock (18\%), jazz (10\%), Latin (5\%), and classical (4\%) genres. BPM range: 60--200. All 24 major/minor keys are represented. Each track was annotated with complete UDMS metadata across all 18 fields by two domain experts; discrepancies were resolved by a third annotator. Software versions: rekordbox 7.1.2, Serato DJ 2.10, Traktor 3.11. Transfers were performed in October 2025 on macOS 14.x.

\subsection{Transfer Protocol}

For each of the six transfer paths ($R \to S$, $R \to T$, $S \to R$, $S \to T$, $T \to R$, $T \to S$), we:
\begin{enumerate}
  \item Exported the full library using the platform's native XML/database export (Rekordbox: XML; Serato: database v1.0; Traktor: NML).
  \item Imported the export into the destination platform using the platform's built-in import functionality.
  \item Re-exported from the destination platform.
  \item Compared field values against the ground-truth UDMS annotation.
\end{enumerate}

This protocol isolates platform-induced degradation from user error. No manual editing was performed between export and import.

\subsection{Preservation Metrics}

We define two metrics. The \textbf{field-level preservation rate} measures per-field fidelity:
\begin{equation}
  \text{PR}(f, p \to q) = \frac{| \{ t \in L : v_p(t, f) = v_q(t, f) \} |}{|L|}
\end{equation}
where $v_p(t, f)$ denotes the value of field $f$ for track $t$ as recorded in platform $p$, compared against the ground-truth UDMS annotation. A field is \emph{preserved} if the transferred value matches the original after type normalization (e.g., BPM rounded to one decimal place is acceptable).

The \textbf{Aggregate Quality Score (AQS)} summarizes transfer fidelity:
\begin{equation}
  \text{AQS}(p \to q) = \frac{1}{|F_{pq}|} \sum_{f \in F_{pq}} w_f \cdot \text{PR}(f, p \to q)
\end{equation}
where $F_{pq}$ is the set of fields present in both $S_p$ and $S_q$, and $w_f$ is a field importance weight (BPM, key, title: $w=2$; all others: $w=1$).

\subsection{UDMS Normalization Pipeline}

The UDMS normalization pipeline converts platform-native metadata to canonical form. Algorithm~\ref{alg:udms-norm} describes the process.

\begin{algorithm}[t]
\caption{UDMS Normalization Pipeline}
\label{alg:udms-norm}
\begin{algorithmic}[1]
\REQUIRE Track metadata $m_p$ from platform $p \in \{R, S, T\}$
\ENSURE Canonical UDMS representation $m_{\text{UDMS}}$
\STATE $m_{\text{UDMS}} \gets \emptyset$
\FOR{each field $f \in F_{\text{UDMS}}$}
  \STATE $v \gets \text{extract}(m_p, f, p)$ \COMMENT{Platform-specific mapping}
  \IF{$f = \text{key}$}
    \STATE $v \gets \text{normalize\_key}(v, p)$ \COMMENT{See Eq.~3}
  \ELSIF{$f = \text{bpm}$}
    \STATE $v \gets \text{round}(v, 1)$ \COMMENT{1 decimal place}
  \ELSIF{$f \in \{\text{genre}, \text{label}\}$}
    \STATE $v \gets \text{normalize\_unicode}(v)$
    \STATE $v \gets \text{title\_case}(v)$
  \ELSIF{$f \in \{\text{energy}, \text{mood}\}$}
    \STATE $v \gets \text{clamp}(v, \text{lo}_f, \text{hi}_f)$
  \ENDIF
  \STATE $m_{\text{UDMS}}[f] \gets v$
\ENDFOR
\RETURN $m_{\text{UDMS}}$
\end{algorithmic}
\end{algorithm}

The key normalization sub-procedure (Eq.~3 in the paper) handles three notation systems: Camelot (e.g., ``8A''), OpenKey (Traktor, e.g., ``4m''), and traditional (e.g., ``Am''). The Python reference implementation maps all variants to Camelot via a lookup table of 24 keys.

\subsection{Degradation Taxonomy}

We classify metadata fields into three tiers based on observed preservation behavior:
\begin{itemize}
  \item \textbf{Tier 1 --- Numeric (audio-derived)}: BPM, bitrate, sample rate, duration. Preserved at $>$99\% across all paths. These values are derived from the audio file itself and are re-computed on import.
  \item \textbf{Tier 2 --- Categorical (human-annotated)}: genre, label, key. Preserve at 66--88\%. Degradation driven by enum mismatches (e.g., ``Electronic'' vs.\ ``EDM''), notation differences (Camelot vs.\ OpenKey), and missing fields.
  \item \textbf{Tier 3 --- Custom (platform-specific)}: energy, mood, rating, playlist structure. Preserve at $<$40\%. Many platforms lack native equivalents; fields are dropped entirely.
\end{itemize}

\subsection{Statistical Analysis}

We use McNemar's test to assess whether preservation rates differ significantly between transfer paths. For each field $f$ and pair of paths $(p \to q, p \to r)$, we construct a $2 \times 2$ contingency table of preserved/degraded counts and compute the McNemar $\chi^2$ statistic. We apply Bonferroni correction for multiple comparisons across all field--path pairs.

\subsection{Auto-Tag Recovery (Experiment~5)}

For tracks with degraded key or genre metadata, we attempt recovery via external sources. We query MusicBrainz and Discogs using ISRC codes (when preserved) or audio fingerprint matching (via AcoustID). Recovery is measured as:
\begin{equation}
  \text{RR}(f) = \frac{| \{ t \in D_f : \text{recovered}(t, f) \} |}{|D_f|}
\end{equation}
where $D_f$ is the set of tracks with degraded field $f$.
